% Part: computability
% Chapter: recursive-functions
% Section: other-recursions

\documentclass[../../../include/open-logic-section]{subfiles}

\begin{document}

\olfileid[hi]{cmp}{rec}{ore}
\olsection{अन्य पुनरावर्तन}

युग्मन और अनुक्रमण का उपयोग करके हम आदिम पुनरावर्तन के अधिक असामान्य
(और उपयोगी) रूपों को उचित ठहरा सकते हैं। मसलन, प्रायः दो फलनों को एक साथ
परिभाषित करना उपयोगी होता है, जैसे निम्न परिभाषा में:
\begin{align*}
h_0(\vec x, 0) & = f_0(\vec x) \\
h_1(\vec x, 0) & = f_1(\vec x) \\
h_0(\vec x, y+1) & = g_0(\vec x, y, h_0(\vec x, y), h_1(\vec x, y)) \\
h_1(\vec x, y+1) & = g_1(\vec x, y, h_0(\vec x, y), h_1(\vec x, y))
\end{align*}
यह \emph{सहकालिक पुनरावर्तन} का एक उदाहरण है। फलनों को परिभाषित करने का
एक अन्य उपयोगी ढंग यह है कि $h(\vec x, y+1)$ का मान \emph{सभी} मानों
$h(\vec x, 0)$, \dots,~$h(\vec x, y)$ के पदों में दिया जाए, जैसे निम्न
परिभाषा में:
\begin{align*}
h(\vec x, 0) & = f(\vec x) \\
h(\vec x, y+1) & = g(\vec x, y, \tuple{h(\vec x, 0), \dots, h(\vec x, y)}).
\end{align*}
निम्न प्रारूप इस विचार को अधिक संक्षेप में व्यक्त करता है:
\[
h(\vec x, y) = g(\vec x, y, \tuple{h(\vec x, 0), \dots, h(\vec x, y-1)})
\]
इस समझ के साथ कि $g$ का अंतिम तर्क, जब $y$, $0$ हो, केवल रिक्त अनुक्रम
है। दोनों निरूपणों में विचार यह है कि ``उत्तराधिकारी चरण'' संगणित करते
समय फलन $h$ अब तक संगणित मानों के पूरे अनुक्रम का उपयोग कर सकता है। इसे
\emph{मान-क्रम पुनरावर्तन} कहते हैं। एक विशिष्ट उदाहरण के रूप में, इसका
उपयोग निम्न प्रकार की परिभाषा को उचित ठहराने के लिए किया जा सकता है:
\begin{align*}
h(\vec x, y) & = \begin{cases}
  g(\vec x, y, h(\vec x, k(\vec x, y))) & \text{यदि $k(\vec x, y) < y$} \\
  f(\vec x) & \text{अन्यथा}
\end{cases}
\end{align*}
दूसरे शब्दों में, $h$ का $y$ पर मान,~$h$ के उस \emph{किसी भी}
पूर्ववर्ती मान पर मान के पदों में संगणित किया जा सकता है जो~$k$ द्वारा दिया गया हो।


\begin{prob}
  मान-क्रम पुनरावर्तन द्वारा शेषफल फलन $r(x,y)$ परिभाषित करें। (यदि
  $x$, $y$ प्राकृतिक संख्याएँ हैं और $y > 0$, तो $r(x,y)$ वह संख्या
  है जो~$y$ से छोटी है तथा $x = z\times y + r(x,y)$ किसी~$z$ के लिए
  मान्य है। निश्चितता के लिए मान लें कि यदि $y=0$, तो $r(x,0) = 0$।)
\end{prob}

आपको विचार करना चाहिए कि सामान्य आदिम पुनरावर्तन का उपयोग करके ये फलन
कैसे प्राप्त किए जाएँ। आदिम पुनरावर्तन का एक अंतिम रूप इस अर्थ में अधिक
लचीला है कि प्रक्रिया के दौरान \emph{प्राचलों} (पार्श्व मानों) को बदलने
की अनुमति होती है:
\begin{align*}
h(\vec x, 0) & = f(\vec x) \\
h(\vec x, y+1) & = g(\vec x, y, h(k(\vec x), y))
\end{align*}
इसका भी सामान्य आदिम पुनरावर्तन से अनुकरण किया जा सकता है। (ऐसा करना
कठिन है। संकेत के लिए संगणन को हाथ से खोलकर देखने का प्रयास करें।)

\end{document}
